% This is the ADASS_template.tex LaTeX file, 24 Oct 2025.
% It is based on the ASP general author template file, but modified to reflect the specific
% requirements of the ADASS proceedings.

% This is the aspauthor.tex LaTeX file
% Copyright 2021, Astronomical Society of the Pacific Conference Series
% Revision:  29 November 2021

% To compile, at the command line positioned at this folder, type:
% pdflatex ADASS_template
% pdflatex ADASS_template
% This will create a file called ADASS_template.pdf

\documentclass[11pt,twoside]{article}

% Do NOT use ANY packages other than asp2023. 
\usepackage{asp2023}

\aspSuppressVolSlug
\resetcounters

% plot articlefiguretwo smaller
%\renewcommand{\plottwo}[2]{%
%  \centering
%  \includegraphics[width=0.4\textwidth]{#1}%
%  \hspace{0.5cm}
%  \includegraphics[width=0.4\textwidth]{#2}%
%}


% References must all use BibTeX entries in a .bibfile.
% References must be cited in the text using \citet{} or \citep{}.
% Do not use \cite{}.
% See ManuscriptInstructions.pdf for more details
\bibliographystyle{asp2021}

% The ``markboth'' line sets up the running heads for the paper.
% 1 author: "Surname"
% 2 authors: "Surname1 and Surname2"
% 3 authors: "Surname1, Surname2, and Surname3"
% >3 authors: "Surname1 et al."
% Replace ``Short Title'' with the actual paper title, shortened if necessary.
% Use mixed case type for the shortened title
% Ensure shortened title does not cause an overfull hbox LaTeX error
% See ASPmanual2010.pdf 2.1.4  and ManuscriptInstructions.pdf for more details
\markboth{Elsa Buchholz, Sagar Borra, Marc Drobek and Hermann He\ss{}ling}{Astronomical Projects with Multiple Repositories}

\begin{document}

\title{Astronomical Projects with Multiple Repositories: A Portal for Consolidated Information and 
Kubernetes-Based Operations}

% Note the position of the comma between the author name and the 
% affiliation number.
% Authors surnames should come after first names or initials, eg John Smith, or J. Smith.
% Author names should be separated by commas.
% The final author should be preceded by "and".
% Affiliations should not be repeated across multiple \affil commands. If several
% authors share an affiliation this should be in a single \affil which can then
% be referenced for several author names. If only one affiliation, no footnotes are needed.
% See ManuscriptInstructions.pdf and ASP's manual2010.pdf 3.1.4 for more details
\author{Buchholz$^1$, Borra$^1$, Drobek$^1$, and  He\ss{}ling$^1$}
\affil{$^1$German Centre for Astrophysics (DZA), G\"{o}rlitz, Saxony, Germany; \email{elsa.buchholz@dzastro.de}}

% This section is for ADS Processing.  There must be one line per author. paperauthor has 9 arguments.
\paperauthor{Elsa Buchholz}{elsa.buchholz@dzastro.de}{}{German Center for Astrophysics (DZA)}{Data Science}{G\"{o}rlitz}{Saxony}{02826}{Germany}
\paperauthor{Sagar Borra}{sagar.borra@dzastro.de}{}{German Center for Astrophysics (DZA)}{Data Science}{G\"{o}rlitz}{Saxony}{02826}{Germany}
\paperauthor{Marc Drobek}{marc.drobek@dzastro.de}{}{German Center for Astrophysics (DZA)}{Data Science}{G\"{o}rlitz}{Saxony}{02826}{Germany}
\paperauthor{Hermann He\ss{}ling}{hermann.hessling@dzastro.de}{}{German Center for Astrophysics (DZA)}{Data Science}{G\"{o}rlitz}{Saxony}{02826}{Germany}


% There should be one \aindex line (commented out) for each author. These are used to
% build up the author index for the Proceedings. The surname must come first, followed by
% initials. Note the use of ~ before each initial to control spacing.
% The \author entries (see above) have surname last. These \aindex entries have
% surname first.
% The Aindex.py command willl create them for you after you have constructed the \author
% The first entry should be the first author, for bold-facing the author index page-reference

%\aindex{Buchholz,~E.}
%\aindex{Borra,~S.}
%\aindex{Drobek,~M.}
%\aindex{He\ss{}ling,~H.}

\begin{abstract}
Machine Learning-based Pipeline for Pulsar Analysis (ML-PPA) aims to enable automated, real-time pulsar searches during data acquisition. Supporting this approach requires not only an efficient analysis pipeline but also a maintainable software infrastructure that facilitates automated documentation, deployment, and centralized access to project information.
To address these challanges, ML-PPA adopts modern software engineering practices, including continuous integration and deployment (CI/CD), automated testing and documentation generation. A central information portal consilidates project information and provides access to documentation to all repositories. The infrastructure is containerised and operated using Kubernetes, enabling automated deployment and updates while keeping software and documentation hosted in-house. The poster provides an overview of the software architecture of the ML-PPA project and explains how the portal and documentation for each of the sub-projects are automatically generated from the source code and made available online. Future developments include Large Languages Models (LLMs) to improve access to the growing body of project documentation.


\end{abstract}

%shorter abstract give in introduction
%science strucutre: 
%What, why, next steps
%The antennas and telescopes of the Square Kilometre Array Observatory (SKAO) will produce enormous data streams, making long-term storage of all raw telescope data impractical. To cope with this amount of data, data must be processed in real time and filtered as close to the telescopes as possible, reducing the volume of data that needs to be stored.


% These lines show examples of subject index entries. At this stage these have to commented
% out, and need to be on separate lines. Eventually, they will be automatically uncommented
% and used to generate entries in the Subject Index at the end of the Proceedings volume.
% Don't leave these in! - replace them with ones relevant to your paper.

    
%\ssindex{algorithm!machine learning}
%\ssindex{software!architecture}
%\ssindex{computing!cluster!Kubernetes}
%\ssindex{packages!Singularity}
%\ssindex{packages!Docker}

% These lines show examples of ASCL index entries. At this stage these have to commented
% out, and need to be on separate lines. Eventually, they will be automatically uncommented
% and used to generate entries in the ASCL Index at the end of the Proceedings volume.
% The ascl.py command will scan your paper on possible code names.
% Don't leave these in! - replace them with ones relevant to your paper.
%\ooindex{FOOBAR, ascl:1101.010}

\section{Introduction}
%0.5 sites
%motivation, problem, goal

The antennas and telescopes of the Square Kilometre Array Observatory (SKAO) will produce enormous data streams\footnote{SKA Observatory, Science Processing Centers, accessed 2025-11-20, URL: \url{https://www.skao.int/en/explore/big-data/361/science-processing-centres}.},
making long-term storage of the raw telescope data impractical. 
Instead, projects such as 'Transients and Pulsars with MeerKAT' (TRAPUM) process the data streams collected per day. They search for pulsars and transients, storing only data of scientific relevance. This processing is performed by a largely automated pipeline, though final steps still require manual inspection \citep{TRAPUM}.
The Machine Learning-based Pipeline for Pulsar Analysis (ML-PPA) project aims to eliminate the manual steps by providing a fully automated analysis pipeline for pulsar searches, developed by an interdisciplinary team of astronomers, physicists and data scientists.
As the ML-PPA project continues to grow, the increasing involvement of researchers and the expansion of software have introduced challenges in maintaining a well-organised and scalable software structure. 

\section{ML-PPA: Concept, Architecture and Software}
%1 site
%ML-PPA: Concept, Architecture and Software
%simulation, rectangle, dynamic filtering to analyse in real time during data aquisistion this needs multi repos software is organised in three layers.
%The overall structure of the current workflow is shown in Fig.~1. The rectangular box shows the simulation pipeline: pulsar signals are generated, the influence of Interstellar Medium (ISM) is taken into account as well as the influence of the noise from the electronics of telescopes.
%The pipelines currently used in radio astronomy are not designed for this purpose. The primary goal of the software package ML-PPA is to provide methods and procedures required for an automated detection of pulsar signals during data acquisition. Machine learning methods for signal detection are presented, and synthetic training data is generated by simulating the path of the signals from the source to detection by the telescopes.  Moreover, real data from the Effelsberg telescope and from MeerKAT are used to check the quality of the ML algorithms and the impact of the electronic noise of telescopes on generating synthetic pulsar signals. 
The workflow of the ML-PPA pipeline has a simulation pipeline included, where pulsar signals are generated. The pipeline includes as well the influence of noise from electronics telescopes and of the Interstallar Medium (ISM). 
Digital Twins are used to train the machine learning algorithms. In addition, real observational data from radio telescopes are used to evaluate the quality of the algorithms. 
% Furthermore, a dynamic filtering mechanism for pulsar signals is included. 
This is intended to enable a search for pulsars during data collection. %and provided directly as data products.
Several neural network architectures have already been analyzed; however, further research is still needed. 

The ML-PPA components are organised in Git repositories. The software itself follows a three layer architecture. The user interface layer provides Jupyter notebooks that allow users to explore the provided software components. The middle layer provides different respositories. Finally, the container level consists of pre-build apptainer containers containing the complete software stack. This ensures that users won't need to install or configure the software environment  in order to run Jupyter notebooks. %Finally, the algorithm layer represents the software componentns for the ML-PPA project.

\articlefiguretwo{P406_f2_a.png}{P406_f2_b.png}{fig2ab}{ ML-PPA System Architecture and Implementations: The goal of the ML-PPA software framework is to provide a fully automated pipeline for pulsar search.  \emph{Left:} The conceptual structure of the current ML-PPA approach, \emph{Right:} the implementation of the project's three-layer software architecture are illustrated in the figures.}

Such project needs to adopt modern engineering practices, including a Continuous Integration and Continuous Deployment (CI/CD) pipeline for each repository, test coverage and automated deployment of documentation. 
In addition to eliminating isolated information silos within the infrastructure, a single portal was set up to centralise all information \citep{ML-PPA_portal}.

\section{ML-PPA: Deployment, Operation and Information Portal}
%1 site
%arc and design desiciion
%Kubernetes Deployment
%ML-PPA, Kubernetes, Container, CICD
% There is a figure command allowing for three figures:
% \articlefigurethree{example.eps}{example.eps}{example.eps}{ex_fig1_triple}{Now there are three of them.}

%\articlefigure[width=0.25\textwidth]{P406_f3.png}{fig3}{Multi-Environment Documentation Deployment:
%Separate CI/CD pipelines for the development and production repositories build documentation from source code, package it into Docker images, and publish them to the Harbor container registry. The images are deployed on separate Kubernetes clusters: an internally accessible development cluster and a publicly accessible production cluster.}



%\articlefigure[width=0.25\textwidth]{P406_f4.png}{fig4}{Fully Automated Documentation Updates:
%Two Git repositories separate tool development from Kubernetes infrastructure management. Keel monitors the documentation images and automatically triggers redeployment when an image update is detected, ensuring that the latest documentation is deployed to the Kubernetes cluster without manual intervention.}
The documentation infrastructure follows the GitOPs-based deployment approach. The automated deployment of documentation is implemented using Kubernetes, Helm, Flux and Keel\footnote{Keel, Documentation, accessed 2026-11-16, URL: \url{https://keel.sh/docs/}.}. Kubernetes provides the infrastructure for hosting the documentation web service. Helm is used to package and manage its deployment in the multi-environment Kubernetes cluster, whie Flux automates the deployment and reconciliation process. For the documentation deployment, Keel replaces Flux  in the Kubernetes cluster to avoid the need for automated Git commits and pushes whenever a new documentation container image becomes available. Instead, Keel monitors the container registry for changes in the image hash (digest) of the documentation container. Whenever a new image is detected, Keel automatically triggers an update of the corresponding Kubernetes deployment. This ensures that the latest version of the documentation is automatically deployed and made available through the web service.
To provide a central access point for all project-related documentation, a dedicated portal was developed to consolidate the available information in one place.


\articlefiguretwo{f3_P406.png}{f4_P406.png}{fig3-4}{  Multi-Environment Documentation Deployment with Fully Automated Documentation Updates \emph{Left:} Separate CI/CD pipelines for the development and production repositories build documentation from source code, package it into Docker images, and publish them to the Harbor container registry. The images are deployed on separate Kubernetes clusters: an internally accessible development cluster and a publicly accessible production cluster. \emph{Right:} Two Git repositories separates tool development from Kubernetes infrastructure management. Keel monitors the documentation images and automatically triggers redeployment when an image update is detected, ensuring that the latest documentation is deployed to the Kubernetes cluster without manual intervention.}

\section{Results and Discussion}
%0.5 sites
%Portal
%review the arc
% setup is a lot of work, documentation  vorteil gegenüber gitlab pages , easy for developers provide docu 
Containerisation and a Kubernetes infrastructure are used to operate the entire ML-PPA system. Both the repositories and the documentation remain in-house, which enables automated updates, thereby reducing maintenance costs in particular. 
The implemented three-layer architecture enables users to access and explore the software through Jupyter notebooks without requiring manual installation of the underlying software stack. Pre-built Apptainer containers ensures a consistent and reproducible execution environment.
Future additions to ML-PPA will follow this workflow to ensure high-quality, well-documented code that is seamlessly deployed on a modern infrastructure.


\articlefigure[width=0.45\textwidth]{P406_f5.png}{fig1}{ML--PPA -- A Central Information Hub: The ML--PPA portal provides a central access point to project information, linking the documentation and GitLab repositories of individual tools. A QR code enables direct access to the portal.}
% It is possible to reduce the size of a figure among other changes (see the instructions).  Here is an example:
% \articlefigure[width=.5\textwidth]{example_image}{ex_fig1_reduced}{Welcome to 1953 a little smaller.}




\section{Conclusion and Outlook}
%0.3 sites
%summary and next steps
Overall, the developed infrastructure brings together the individual software components into a unified and accessible environment for the ML-PPA project. The infrastructure enables the reproducible and integrated development of machine learning algorithms for pulsar search.
The already substantial volume of software documentation makes it difficult to extract relevant information. Therefore, LLMs (Large Language Models) are to be integrated in future to facilitate the use of ML-PPA. One possible approach is to extend existing models using a skill-based approach. Skills provide reusable, task-specific instructions and contextual information that guide LLMs in performing specilasied tasks without requiring modifications to the underlying models \citet{agent_skills_specification}. This can help users to find the correct parameters and configuration of the software in addition to the documentation.

\acknowledgements 
    %We would like to thank the Center for Information Services and High Performance Computing (ZIH) of TU Dresden for providing computational resources and support. 
    The authors gratefully acknowledge the financial support provided by the German Federal Ministry of Research, Technology and Space (BMFTR) in the framework of the Knowledge creates perspectives for the region, for the project StStG -- DZA -- Aufbauphase: Deutsches Zentrum f\"{u}r Astrophysik, Gro\ss{}forschungszentrum in der s\"{a}chsischen Lausitz: Aufbauphase 2026, grant number 03WSP1746."


\bibliography{P406}  % For BibTex

% if we have space left, we might add a conference photograph here. Leave commented for now.
% \bookpartphoto[width=1.0\textwidth]{foobar.eps}{FooBar Photo (Photo: Any Photographer)}

\end{document}
